\item Considere una película delgada de índice de refracción $n$ de espesor $t$ rodeada por el vacío. Un rayo de luz incide sobre la película con un ángulo de incidencia oblicuo $\theta_1$ con respecto a la normal. Encuentre la relación matemática general que debe cumplir el espesor $t$, el índice $n$ de la película, el ángulo de incidencia $\theta_1$ y la longitud de onda de la luz $\lambda$ para que ocurra interferencia constructiva en los rayos reflejados.